% Generated draft: requires mathematical and source review.
\documentclass[11pt]{article}
\usepackage[margin=25mm]{geometry}
\usepackage{fontspec}
\setmainfont{Latin Modern Roman}
\usepackage{amsmath,amssymb,mathtools,unicode-math}
\setmathfont{Latin Modern Math}
\usepackage{enumitem,needspace,xurl,hyperref}
\setlength{\parindent}{0pt}
\setlength{\parskip}{5pt}
\setlength{\emergencystretch}{3em}
\newcommand{\N}{\mathbb{N}}
\newcommand{\Z}{\mathbb{Z}}
\newcommand{\Q}{\mathbb{Q}}
\newcommand{\R}{\mathbb{R}}
\newcommand{\C}{\mathbb{C}}
\newcommand{\F}{\mathbb{F}}
\newcommand{\A}{\mathbb{A}}
\newcommand{\PP}{\mathbb{P}}
\newcommand{\E}{\mathbb{E}}
\newcommand{\eps}{\varepsilon}
\newcommand{\dd}{\,\mathrm{d}}
\begin{document}
\section*{UnsolvedMath ID: 20003290}
\textbf{Source problem code:} catalogue entry
\subsection{Definitions and mathematical statement}
The supplied verification identifies this permanent record as explanatory text between Conjecture 20 (Rudnick--Sarnak) and Conjecture 21, not an independent conjecture. Neither neighboring conjecture is substituted. Its setting is a finite-area hyperbolic surface $X=\Gamma\backslash\mathbb H^2$ with cusp width one, meaning $\Gamma\cap U=\{n_m:m\in\Z\}$ for $n_t=\begin{pmatrix}1&t\\0&1\end{pmatrix}$. In $Y=\Gamma\backslash PSL(2,\R)\cong T^1X$ put $a_y=\operatorname{diag}(\sqrt y,1/\sqrt y)$ and $u_y(t)=\Gamma n_ta_y$. Use the surface convention, or the corresponding orbifold version if torsion is allowed. Then $u_y$ has parameter period one and physical horocycle length $y^{-1}$.

The captured context reports that the normalized continuous horocycle measure tends to Liouville probability $m_Y$ as $y\to0^+$, while for fixed $y$ and irrational $\alpha$, $u_y(n\alpha)$ samples that circle by irrational rotation. The samples use $n\alpha$, not $\alpha^n$. These two separate limits do not specify a simultaneous relation between sample count and height.

The aid's quantitative companion fixes a smooth compactly supported observable $f$ and writes $F_y(t)=f(u_y(t))$. Suppose
\[
\|k\alpha\|_{\R/\Z}\geq c|k|^{-\tau}\quad(k\in\Z\setminus\{0\}),\qquad c>0,
\]
and let the integer $s>\tau+1$. The supplied small-divisor estimate is
\[
\left|\frac1M\sum_{n=1}^MF_y(n\alpha)-\int_0^1F_y(t)\,\dd t\right|
\leq\frac{\zeta(s-\tau)}{c(2\pi)^sM}\,\|F_y^{(s)}\|_{L^1([0,1])}.
\]
Here the Riemann zeta function comes from summing Fourier modes; no Diophantine bound is being assumed for every irrational $\alpha$. A complementary estimate for bounded variation is
\[
\left|\frac1M\sum_{n=1}^MF_y(n\alpha)-\int_0^1F_y\right|
\leq C_\tau\operatorname{Var}(F_y)\,c^{-1/(\tau+1)}M^{-1/(\tau+1)}.
\]
The constant $C_\tau$ is independent of $M,y,F_y$ under these hypotheses.

Let $Hf(g)=\frac{\dd}{\dd r}|_{r=0}f(gn_r)$ be the physical right-horocycle derivative. The exact scaling is
\[
F_y^{(s)}(t)=y^{-s}(H^sf)(u_y(t)),\qquad
\|F_y^{(s)}\|_1=y^{-s}\int_0^1|(H^sf)(u_y(t))|\,\dd t.
\]
Defining $E_{\mathrm{hor}}(y,f)=|\int_0^1F_y-\int_Yf\,\dd m_Y|$, the total discrete-to-Liouville error is at most the Fourier bound plus $E_{\mathrm{hor}}(y,f)$, or the bounded-variation bound plus that term. The context supplies $E_{\mathrm{hor}}(y,f)\to0$ but no rate.

At the critical proposed scale $y\asymp M^{-2}$, the elementary bound $\|F_y^{(s)}\|_1\leq y^{-s}\|H^sf\|_\infty$ makes the first error estimate grow like $M^{2s-1}$ rather than tend to zero. This shows insufficiency of that estimate and of separate qualitative limits; it does not prove failure of the actual simultaneous equidistribution. No missing Conjecture 21 or critical-scale theorem is asserted by this context-only record.
\subsection{Short English statement}
This is the paragraph between two conjectures. Its supplied companion bounds separate fixed-horocycle sampling error from continuous horocycle equidistribution error, and displays the height-dependent derivatives explicitly. Those derivatives make the elementary bound ineffective at height comparable to the inverse square of sample count. This is a quantitative gap, not a disproof of simultaneous equidistribution or a replacement conjecture.
\subsection{Sources}
\begin{itemize}
\item Local registry: permanent catalogue entry, exact source code catalogue entry; \texttt{aim-other-notes.json}, \texttt{notes[96]}; inherited canonical tag ``conjecture''.
\item AIM \emph{Open Problems from the Workshop ``Emerging Applications of Measure Rigidity''}, Section 4.4, \url{https://aimath.org/WWN/measrigid/measrigid.pdf}. The supplied audit places this paragraph after Rudnick--Sarnak Conjecture 20 and before Conjecture 21. No independent numbered conjecture or proposer is present.
\item Supplied research attempt 1, registry updated 14 August 2026: \texttt{context\_only}, \texttt{not\_a\_problem}, low novelty confidence, machine-generated and not peer reviewed. It supplies the Fourier constant, BV estimate, physical derivative scaling, and two-scale ledger, without a complete authoritative reference for the continuous-horocycle theorem.
\item Statement verification reports the exact contextual boundary in the official PDF; registry status nevertheless remains open/OPEN-TRIAGE. Original remarks and literature field are empty.
\end{itemize}
\end{document}
